% !TeX program = xelatex
% !TeX encoding = UTF-8 Unicode
\documentclass[UTF8,12pt,a4paper]{ctexart}

\usepackage[a4paper,top=2.4cm,bottom=2.4cm,left=2.6cm,right=2.6cm]{geometry}
\usepackage{booktabs}
\usepackage{array}
\usepackage{tabularx}
\usepackage{caption}
\usepackage{xcolor}
\usepackage{enumitem}
\usepackage{hyperref}

\hypersetup{hidelinks}
\setlength{\parskip}{4pt}
\linespread{1.30}

\captionsetup{font=small,labelfont=bf}
\renewcommand{\arraystretch}{1.18}

\title{基于 yolo11 n/s/m 三模型 HN 全量比例扫描的对比分析}
\author{}
\date{2026 年 4 月 11 日}

\begin{document}
\maketitle

\section*{0. 数据来源}
\addcontentsline{toc}{section}{0. 数据来源}

本分析基于以下三组在统一第一阶段门控协议下完成的 HN 全量比例扫描资产,排序规则统一为 \texttt{Spec@R99.5} 降序、\texttt{Spec@R99.0} 降序、\texttt{Prec@R99.0} 降序、\texttt{PTR@R99.0} 升序;评估均来自校准集 (val-cal, 216 张) + 工作点评估集 (val-op, 504 张, 含 84 个 normal)。

\begin{itemize}[leftmargin=*,itemsep=2pt]
    \item \textbf{yolo11n-cls} sweep:\texttt{research/results/stage1\_formal/gate\_hn\_n\_sweep/}
    \item \textbf{yolo11s-cls} sweep:\texttt{research/results/stage1\_formal/gate\_hn\_s\_sweep/} (本次新增)
    \item \textbf{yolo11m-cls} sweep:\texttt{research/results/stage1\_formal/gate\_hn\_m\_sweep/}
\end{itemize}

\noindent 三组各 11 档比例 (\texttt{hn00, hn02, hn04, \ldots, hn20}),共 33 个数据点;hn00 均为容量扫描基线复用。所有 \(\Delta\) 与排名以下文表格为准。

\bigskip
\section*{1. HN 收益普遍存在,但最优 ratio 完全不可迁移}
\addcontentsline{toc}{section}{1. HN 收益普遍存在但最优 ratio 完全不可迁移}

\begin{table}[h]
\centering
\caption{三模型 HN 最优比例点对照}
\begin{tabular}{lcccccc}
\toprule
模型 & 参数量 & hn00 baseline & 最优 ratio & 最优 spec995 & \(\Delta\)spec995 & 多判对 normal \\
\midrule
yolo11n-cls & 5.4M  & 0.5119 & \textbf{hn08} & 0.5833 & +0.0714 & +6 \\
yolo11s-cls & 10M   & 0.5238 & \textbf{hn20} & 0.5952 & +0.0714 & +6 \\
yolo11m-cls & 20.1M & 0.5952 & \textbf{hn14} & 0.6429 & +0.0476 & +4 \\
\bottomrule
\end{tabular}
\end{table}

\textbf{结论}:三个容量模型的最优比例点分别落在 \texttt{hn08} / \texttt{hn14} / \texttt{hn20},且容量与最优 ratio 之间不存在任何单调或线性关系。这把 "HN 不存在普适最优 ratio" 的判断从原来 m+x 两点的弱证据加强到 \textbf{n+s+m 三点的强证据 —— HN ratio 是骨干网络强相关超参数,不能跨模型移植}。三模型只有 "存在最优点" 这件事是共同的;最优点在哪里完全是模型属性。

\bigskip
\section*{2. 三模型非单调形态对比 \mdseries (低比例启动效应在 s 上不成立)}
\addcontentsline{toc}{section}{2. 三模型非单调形态对比}

按 ratio 升序观察 \texttt{Spec@R99.5} 的形态:

\begin{table}[h]
\centering
\caption{三模型 spec995 随 ratio 的变化形态(粗体为各自最优)}
\small
\begin{tabular}{lccccccccccc}
\toprule
ratio & hn00 & hn02 & hn04 & hn06 & hn08 & hn10 & hn12 & hn14 & hn16 & hn18 & hn20 \\
\midrule
n & 0.5119 & 0.5595 & 0.5238 & 0.5357 & \textbf{0.5833} & 0.5476 & 0.5476 & 0.5238 & 0.5476 & 0.5714 & 0.5595 \\
s & 0.5238 & 0.5119 & 0.5595 & 0.5000 & 0.5357 & 0.5833 & 0.5238 & 0.5357 & 0.5476 & 0.5476 & \textbf{0.5952} \\
m & 0.5952 & 0.6429 & 0.5595 & 0.5952 & 0.5833 & 0.5238 & 0.5357 & \textbf{0.6429} & 0.5714 & 0.5476 & 0.5595 \\
\bottomrule
\end{tabular}
\end{table}

形态特征:
\begin{itemize}[leftmargin=*,itemsep=2pt]
    \item \textbf{n}:双峰 (\texttt{hn02} 起步 + \texttt{hn08} 中段),后段 hn14 跌回基线水平
    \item \textbf{s}:后置单峰 + 中段塌陷,\texttt{hn06} 跌至 \textbf{0.5000} 全表谷底,直到 \texttt{hn20} 才达到 0.5952
    \item \textbf{m}:双峰 (\texttt{hn02} 起步 + \texttt{hn14} 中段),hn10 跌至 0.5238 谷底
\end{itemize}

\textbf{尤其值得注意的低比例启动效应不一致}:\texttt{hn00} \(\rightarrow\) \texttt{hn02} 这一步,n 上升 +0.0476、m 上升 +0.0476,但 \textbf{s 反而下降 -0.0119},直接落入 11 档中的次差档位 (rank 10)。\textbf{"低比例 HN 启动收益" 在 m/n 上稳健存在,在 s 上被打破}。这是非单调性中的一个具体的模型异质性证据。

\bigskip
\section*{3. 「充分回流追赶效应」成立但有清晰边界}
\addcontentsline{toc}{section}{3. 充分回流追赶效应及边界}

最关键的对照:\texttt{s+hn20} \(\leftrightarrow\) \texttt{m+hn00}

\begin{table}[h]
\centering
\caption{s 充分回流后与 m 基线的全 4 指标对比}
\begin{tabular}{lcccc}
\toprule
设置 & Spec@R99.5 & Spec@R99.0 & Prec@R99.0 & PTR@R99.0 \\
\midrule
m+hn00 (主线基线) & 0.5952 & 0.6548 & 0.9348 & 0.8829 \\
s+hn20 (s 当前最优) & 0.5952 & 0.6071 & 0.9267 & 0.8929 \\
\midrule
差值 (s \(-\) m) & \textbf{0.0000} & \textbf{-0.0476} & \textbf{-0.0081} & \textbf{+0.0100} \\
\bottomrule
\end{tabular}
\end{table}

s 充分回流后在主排序指标 \texttt{Spec@R99.5} 上 \textbf{精确追平} m 基线 (双方都在 84 个 normal 中判对 50 个),但在 \texttt{Spec@R99.0} 仍低 4 个正常样本,\texttt{Prec@R99.0} 也明显低,通过率 PTR 反而更高。

\textbf{结论}:HN 只能把 s 的 "高召回拦截" 推到 m 基线水平,\textbf{无法消除容量本身带来的次级指标差距}。换言之 "充分回流的较小容量可以替代未回流的主模型" 这句话\textbf{只在主排序指标上成立,在次级约束上不成立}。

\bigskip
\noindent 更严格的检验:整个 33 个数据点中,有没有任何一组 (s 或 n 的) 设置能超过 \texttt{m+hn14} 的全局最强 (spec995=0.6429) ?

\textbf{答案是没有}。最高的两个数据点是 \texttt{m+hn14} 与 \texttt{m+hn02} (均为 0.6429),其后并列第三的是 \texttt{m+hn00}、\texttt{m+hn06} 与 \texttt{s+hn20} (均为 0.5952);全 33 个数据点的全局最强始终来自 m 模型自身。这给论文第 7 章关于 \textbf{"HN 不能替代主模型选择"} 的工程结论提供了直接的反例搜索证据。

\bigskip
\section*{4. 提升幅度与基线强度负相关 \mdseries (基线越弱,HN 抬升越大)}
\addcontentsline{toc}{section}{4. 提升幅度与容量负相关}

n 和 s 的最优 \(\Delta\)spec995 均为 \textbf{+0.0714} (恰好都是多 6 个正常样本判对),而 m 的最优 \(\Delta\)spec995 仅为 \textbf{+0.0476} (4 个正常样本)。HN 对 \textbf{基线较弱} 的较小容量模型抬升幅度明显更大。

机制解读:小容量模型在 normal 尾部分布建模上更欠缺,HN 有更多的"边界纠错"空间;而 m 在基线时已经较强,HN 提供的额外信号边际效用递减。但要特别注意:这\textbf{不是} "小容量模型用 HN 后更强",而是 \textbf{"小容量模型用 HN 后从更低位置抬升更多 —— 但终点仍未追上 m+hn14"}。"HN 抬升幅度大" 不等于 "HN 后绝对水平高"。

\bigskip
\section*{5. \texttt{best\_ep=1} 现象与容量负相关 \mdseries (容量越小越早过拟合)}
\addcontentsline{toc}{section}{5. best ep=1 现象与容量负相关}

\begin{table}[h]
\centering
\caption{各模型 sweep 中 best epoch=1 (第一轮即触顶) 的比例数}
\begin{tabular}{lcccc}
\toprule
模型 & best\_ep=1 的比例数 / 11 & 出现位置 & 占比 \\
\midrule
yolo11n & 6 & hn06, hn10, hn12, hn16, hn18, hn20 & 55\% \\
yolo11s & 4 & hn08, hn12, hn16, hn18 & 36\% \\
yolo11m & 2 & hn08, hn20 & 18\% \\
\bottomrule
\end{tabular}
\end{table}

容量越小,越多 ratio 在第 1 个 epoch 就触顶,之后训练只在退化。这与 ImageNet 预训练初值在小模型上 "几乎已经够用" 一致 —— 小模型表达能力有限,容易在 HN 副本数较高时立即陷入过拟合;m 模型则有足够容量去跟上 HN 信号,只在 hn08/hn20 这两档(显著偏离最优 ratio)才出现 best\_ep=1 的早熟。

\textbf{这一观察现在有 n/s/m 三点支撑,可以在第 5 章训练动力学讨论里直接作为证据引用 —— 它把 "容量小则训练动力学更不稳定" 从一个直觉判断转为可量化的统计现象 (55\% / 36\% / 18\%)}。

\bigskip
\section*{6. 全 4 指标对称改善的最优点只有一个: s+hn20}
\addcontentsline{toc}{section}{6. 全 4 指标对称改善只有 s hn20}

把每个模型的 "最优行" 跟 "自己的 hn00" 比 4 个指标的方向:

\begin{table}[h]
\centering
\caption{三模型最优 ratio 相对各自 hn00 的 4 指标 \(\Delta\)}
\begin{tabular}{lcccccc}
\toprule
模型 + 最优 ratio & \(\Delta\)spec995 & \(\Delta\)spec990 & \(\Delta\)prec990 & \(\Delta\)ptr990 & 全 4 指标改善 \\
\midrule
n+hn08 & +0.0714 & +0.0000 & +0.0003 & \textcolor{red}{+0.0040} & 否 (ptr 恶化) \\
\textbf{s+hn20} & \textbf{+0.0714} & \textbf{+0.0714} & \textbf{+0.0124} & \textbf{-0.0099} & \textbf{是} \\
m+hn14 & +0.0476 & +0.0000 & +0.0001 & \textcolor{red}{+0.0020} & 否 (ptr 恶化) \\
\bottomrule
\end{tabular}
\end{table}

\textbf{只有 s+hn20 在 4 个指标上同时改善}。n+hn08 和 m+hn14 都伴随 ptr990 的轻微上升 (通过率反而升,意味着拦截绝对量没相应放大,只是工作阈值发生了移动)。

这一发现把原论文 "HN 收益不是所有指标都同步改善" 的判断进一步细化为 \textbf{"在 m 和 n 上都不是同步,只有 s 是个例外"}。s 的"质量纯度"是当前 sweep 中 HN 最干净的一个数据点,虽然其绝对水平不及 m+hn14。这反向提示:s 上的 HN 训练可能是后续高价值困难样本评估实验最值得作为切换教师模型的候选。

\bigskip
\section*{一句话总括}

\noindent\textbf{当前 n/s/m 三组 sweep 共同说明}:HN 的 "是否有用" 已经稳健成立 (3/3 模型有正向最优点),但 HN 的 "最优 ratio" 完全不可迁移 (n=8\%, s=20\%, m=14\%)、提升幅度与基线强度呈反向 (越弱抬升越多)、且只有 s+hn20 是 4 指标对称改善的最优点;最重要的是 —— \textbf{HN 无法跨越容量阶梯,主模型 m+hn14 仍是 33 个数据点中的全局最强 (spec995=0.6429)}。这把第 5 章原来 m+x 两点支撑的"模型相关"判断升级为 n+s+m 三点支撑的 \textbf{"具有具体形态特征的模型相关性"},而第 7 章 "HN 不能替代主模型选择" 的工程结论也得到了直接的反例搜索验证 (33 个数据点全局最强仍来自 m)。

\bigskip
\section*{统计学边界声明}

工作点评估集只有 84 个 normal,故 spec995 的最小步长为 \(1/84 \approx 0.0119\)。本文所有 \(\Delta\)spec995 提升的绝对量级在 4--6 个正常样本之间,且当前未补做多随机种子重复或 bootstrap 区间。因此上述全部结论应理解为 \textbf{当前数据划分与当前随机种子下的点估计排序},而非已获得统计显著性的稳定差异。后续若要把任一结论写成 "稳定成立" 的判断,仍需补多 seed 或 bootstrap 验证。

\end{document}
